\documentclass[12pt,a4paper]{article}
\usepackage{fontspec}
\setmainfont{DejaVu Serif}
\setsansfont{DejaVu Sans}
\setmonofont{DejaVu Sans Mono}
\usepackage{geometry}
\usepackage{hyperref}
\usepackage{booktabs}
\usepackage{longtable}
\usepackage{array}
\usepackage{enumitem}
\usepackage{xcolor}
\usepackage{titlesec}
\usepackage{setspace}

\geometry{a4paper, left=24mm, right=24mm, top=28mm, bottom=28mm}
\onehalfspacing

\definecolor{keywordblue}{rgb}{0.0,0.2,0.7}
\titleformat{\section}{\large\bfseries\color{keywordblue}}{\thesection}{1em}{}
\newcolumntype{P}[1]{>{\raggedright\arraybackslash}p{#1}}

\title{CodeAlchemist -- 8. Hafta Kod Bazli Teknik Dokuman}
\author{CodeAlchemist Team}
\date{Mart 2026}

\begin{document}
\maketitle
\tableofcontents
\newpage

\section{Kapsam}
Bu dokuman, 8. haftada hayata alinan 4 ozelligin kod seviyesindeki implementasyonunu aciklar:

\begin{enumerate}[leftmargin=*]
  \item Baglam optimizasyon asistani
  \item Patch guvenlik seviyesi
  \item Paylasimda canli isbirligi (yorum + onay/revizyon dongusu)
  \item Free/Premium plan ayrimi (istek/token kotasi)
\end{enumerate}

Hedef, urun gereksinimlerini dogrudan fonksiyonlar, endpointler ve UI davranisi ile eslestirmektir.

\section{Hafta 7'den Hafta 8'e Donusum (Once / Simdi)}
Bu bolum, 8. haftadaki gelistirmelerin neden degerli oldugunu gostermek icin
"once nasildi, simdi nasil calisiyor" perspektifini verir.

\begin{longtable}{@{} P{3.6cm} P{4.8cm} P{5.8cm} @{} }
\toprule
Alan & Once (Hafta 6-7) & Simdi (Hafta 8) \\
\midrule
Baglam yonetimi & Context bar bilgi veriyordu, kullanici elle duzenliyordu & Prompt oncesi optimize onerisi ve auto optimize akisi var \\
Patch uygulama & Patch dogrudan uygulanabiliyordu & Risk skoru ile safe/review/block karari devreye girdi \\
Paylasim & Link ile paylasim ve goruntuleme odakliydi & Review status + yorum + onay/revizyon + eszamanli iterasyon eklendi \\
Planlama / maliyet & Tuketim gorunurlugu sinirliydi & Free/Premium plan, gunluk istek ve aylik token kotasi aktif \\
\bottomrule
\end{longtable}

\section{Mimari Ozet}
\subsection{Frontend}
\begin{itemize}[leftmargin=*]
  \item Ana orkestrasyon: \texttt{client/src/App.jsx}
  \item Chat ve patch akisi: \texttt{client/src/components/ChatInterface.jsx}
\end{itemize}

\subsection{Backend}
\begin{itemize}[leftmargin=*]
  \item API ve is kurallari: \texttt{server/app.py}
  \item Veri modelleri: \texttt{server/models.py}
\end{itemize}

\section{Ozellik 1: Baglam Optimizasyon Asistani}
\subsection{Urun Gereksinimi}
Prompt gondermeden once, kaliteyi koruyarak gereksiz token yukunu azaltacak oneriler sunulmasi.
Auto mode acik oldugunda sistemin bunu akis icinde otomatik uygulayabilmesi.

\subsection{Kod Entegrasyonu}
\begin{itemize}[leftmargin=*]
  \item \texttt{client/src/components/ChatInterface.jsx:194}
  \begin{itemize}[leftmargin=*]
    \item \texttt{optimizePromptWithReport(prompt)}
    \item Prompt icin before/after token tahmini ve iyilestirme raporu uretir.
  \end{itemize}

  \item \texttt{client/src/components/ChatInterface.jsx:876}
  \begin{itemize}[leftmargin=*]
    \item Gonderim oncesi optimize raporunu devreye alir.
  \end{itemize}

  \item \texttt{client/src/components/ChatInterface.jsx:1527}
  \begin{itemize}[leftmargin=*]
    \item \texttt{Auto Optimize: ON/OFF} UI kontrolu.
  \end{itemize}
\end{itemize}

\subsection{Davranis}
\begin{enumerate}[leftmargin=*]
  \item Kullanici prompt yazar.
  \item Auto optimize aciksa rapor olusturulur.
  \item Kullanici optimize metni uygular, optimize haliyle gonderir veya orijinal metinle devam eder.
\end{enumerate}

\subsection{Kisa Senaryo (Once / Simdi)}
\begin{itemize}[leftmargin=*]
  \item \textbf{Once:} Kullanici 2-3 farkli kod parcasi ve uzun aciklama ile prompt gonderiyor,
  token yukunu ancak cevap gec gelince fark ediyordu.
  \item \textbf{Simdi:} Gondermeden hemen once sistem tekrarli kisimlari sikistirip
  tahmini token kazancini gosteriyor; kullanici optimize veya orijinal metinle devam etmeyi seciyor.
\end{itemize}

\subsection{Tekrar Disi Optimizasyon Onerileri ve Ornek Hatali Promptlar}
Asagidaki maddeler, sadece tekrar tespiti disinda kalan baglam optimizasyon alanlaridir.
Her madde icin bir adet ``hatali prompt'' ornegi verilmis, sistemin onermesi beklenen iyilestirme de belirtilmistir.

\begin{longtable}{@{} P{0.6cm} P{4.0cm} P{5.2cm} P{3.6cm} @{} }
\toprule
No & Optimizasyon fikri & Ornek hatali prompt & Beklenen onerme \ 
\midrule
1 & Gereksiz baglam budama &
``Bugun kahve ictim, toplanti yaptim, sonra markete gittim, bu arada kodu da duzelt.'' &
Teknik hedefi tut, hikaye/gurultuyu kaldir \ 

2 & Amac netlestirme &
``Bunu hizlandir, guvenli yap, UI tasarla, deploy et, test yaz.'' &
Tek ana hedef sec ve onceliklendir \ 

3 & Cikti formati standardizasyonu &
``Nasil uygun gorursen oyle cevapla.'' &
Tablo/madde/kod blogu gibi net format belirt \ 

4 & Kisit ekleme &
``Bu problemi cozer misin?'' &
Dil, uzunluk, kapsam, karmaşıklik siniri ekle \ 

5 & Rol ve hedef kitle uyumu &
``Bu kodu anlat.'' &
Kime anlattigini belirt (junior/senior/PM) \ 

6 & Kod sinyali guclendirme &
``Asagidakini iyilestir: def f(x): return x*x'' &
Dil ve istenen cikti tipini netlestir (diff/fonksiyon) \ 

7 & Test odakli tamamlama &
``Bu fonksiyonu duzelt.'' &
Edge-case ve basit test maddeleri ekle \ 

8 & Risk ve guvenlik filtresi &
``Auth token akisini komple degistir, direkt patch ver.'' &
Risk analizi + rollback plani ile ilerle \ 

9 & Token maliyeti azaltma &
``1200 satir logu komple analiz et.'' &
Yalniz ilgili log kesitlerini iste \ 

10 & Cok adimli gorev ayrisimi &
``Bu urunu bastan sona optimize et.'' &
Analiz/cozum/test adimlarina bol \ 

11 & Belirsizlik tespiti &
``Calismiyor, duzelt.'' &
Hata mesaji, tekrar adimi, ortam bilgisi iste \ 

12 & Dil kalitesi sadeleme &
``Ayni seyi tekrar ediyorum, iyi olsun, daha iyi olsun.'' &
Tekrari azalt, net ve kisa hedef cumlesi uret \ 
\bottomrule
\end{longtable}

\subsection{UI Tarafinda Beklenen Davranis (Ozellik-1 icin)}
Bu 12 durumdan biri algilandiginda, optimize karti asagidaki aksiyonlardan en az birini sunmalidir:
\begin{itemize}[leftmargin=*]
  \item \texttt{Apply Optimized Text}
  \item \texttt{Send Optimized}
  \item \texttt{Send Original}
\end{itemize}

\section{Ozellik 2: Patch Guvenlik Seviyesi}
\subsection{Urun Gereksinimi}
Patch uygulamadan once risk skorlamasi yapilarak 3 modda karar verilmesi:
\texttt{guvenli uygula}, \texttt{incele}, \texttt{blokla}.

\subsection{Kod Entegrasyonu}
\begin{itemize}[leftmargin=*]
  \item \texttt{client/src/components/ChatInterface.jsx:241}
  \begin{itemize}[leftmargin=*]
    \item \texttt{evaluatePatchRisk(currentCode, incomingCode)}
    \item Risk skoru ve gerekce listesi hesaplar.
  \end{itemize}

  \item \texttt{client/src/components/ChatInterface.jsx:436}
  \begin{itemize}[leftmargin=*]
    \item Patch uygulama adimi risk degerlendirmesinden gecirilir.
  \end{itemize}

  \item \texttt{client/src/components/ChatInterface.jsx:2015}
  \begin{itemize}[leftmargin=*]
    \item \texttt{Patch Safety Check} modal gorunumu.
  \end{itemize}
\end{itemize}

\subsection{Risk Faktorleri}
\begin{itemize}[leftmargin=*]
  \item Degisim buyuklugu
  \item Silinen satir orani
  \item Kritik anahtar kelimeler (or. auth/token/secret/billing)
  \item Test etkisi sinyali
\end{itemize}

\subsection{Davranis}
\begin{itemize}[leftmargin=*]
  \item Dusuk risk: guvenli uygulama
  \item Orta risk: inceleme modalinden kontrollu devam
  \item Yuksek risk: otomatik bloklama
\end{itemize}

\subsection{Kisa Senaryo (Once / Simdi)}
\begin{itemize}[leftmargin=*]
  \item \textbf{Once:} Buyuk bir patch tek tikla uygulanabiliyor,
  kritik dosya etkisi sonradan fark ediliyordu.
  \item \textbf{Simdi:} Uygulama oncesi risk skoru cikiyor; sistem
  guvenli ise uygula, orta riskte incelet, yuksek riskte blokla kararini veriyor.
\end{itemize}

\section{Ozellik 3: Paylasimda Canli Isbirligi Modu}
\subsection{Urun Gereksinimi}
Sadece paylasim linki yerine takim ici inceleme dongusu:
yorum yazma, durum degistirme, onay veya revizyon talebi.

Bu haftadaki fark: ekip uyeleri ayni paylasim oturumu uzerinde
yorumlasirken, patch/revizyon kararlarini da ayni akis icinde yurutebiliyor.
Yani "once sadece link" varken, simdi "link + yorum + durum + revizyon" birlikte calisiyor.

\subsection{Veri Modeli}
\begin{itemize}[leftmargin=*]
  \item \texttt{server/models.py:72} \texttt{CollaborationReview}
  \item \texttt{server/models.py:85} \texttt{CollaborationComment}
\end{itemize}

\subsection{API Entegrasyonu}
\begin{itemize}[leftmargin=*]
  \item \texttt{server/app.py:6195} \texttt{GET /api/collaboration/session/<token>/review}
  \item \texttt{server/app.py:6221} \texttt{POST /api/collaboration/session/<token>/review/comment}
  \item \texttt{server/app.py:6256} \texttt{POST /api/collaboration/session/<token>/review/status}
\end{itemize}

\subsection{Frontend Entegrasyonu}
\begin{itemize}[leftmargin=*]
  \item \texttt{client/src/App.jsx:254} \texttt{fetchCollabReview}
  \item \texttt{client/src/App.jsx:269} \texttt{submitCollabReviewComment}
  \item \texttt{client/src/App.jsx:293} \texttt{updateCollabReviewStatus}
  \item \texttt{client/src/App.jsx:2560-2572} durum butonlari
\end{itemize}

\subsection{Durum Makinesi}
\begin{itemize}[leftmargin=*]
  \item \texttt{open}
  \item \texttt{revision\_requested}
  \item \texttt{approved}
\end{itemize}

\section{Ozellik 4: Free/Premium Plan Ayrimi}
\subsection{Urun Gereksinimi}
Kullanicilarin plan bazli istek ve token limitleriyle yonetilmesi,
limit asiminda premiuma gecis mesaji verilmesi.

\subsection{Once / Simdi}
\begin{itemize}[leftmargin=*]
  \item \textbf{Once:} Kullanicilarin maliyet kontrolu dogrudan urun akisi icinde net degildi.
  \item \textbf{Simdi:} Her istek oncesi plan-kota kontrolu yapiliyor; limitte sistem
  kontrollu olarak istegi reddedip premiuma gecis mesaji donuyor.
\end{itemize}

\subsection{Plan Modeli ve Limitler}
\texttt{PLAN\_LIMITS} sabiti ile iki plan tanimli:
\begin{itemize}[leftmargin=*]
  \item \textbf{Free}
  \begin{itemize}[leftmargin=*]
    \item \texttt{daily\_requests = 30}
    \item \texttt{monthly\_tokens = 120000}
  \end{itemize}
  \item \textbf{Premium}
  \begin{itemize}[leftmargin=*]
    \item \texttt{daily\_requests = 250}
    \item \texttt{monthly\_tokens = 1200000}
  \end{itemize}
\end{itemize}

Bu model, hem kisa vadeli suistimali (gunluk istek) hem de uzun vadeli maliyeti
(aylik token) ayni anda kontrol eder.

\subsection{Kota Muhasebesi Nasil Calisiyor?}
\begin{enumerate}[leftmargin=*]
  \item Sistem kullanicinin \texttt{preferences} alanindan plan bilgisini okur.
  \item Gunluk saya\c{c} (tarih bazli) ve aylik token sayaci (ay bazli) init/reset edilir.
  \item Istek gelince tahmini token hesaplanir:
  \begin{itemize}[leftmargin=*]
    \item yaklasik: toplam karakter / 4
  \end{itemize}
  \item \texttt{consume\_plan\_quota(...)} fonksiyonu ile
  projected degerler limitlerle karsilastirilir.
  \item Limit asilmazsa saya\c{c}lar commit edilir ve istek devam eder.
  \item Limit asilirse 429 donulur, sebep kodu ve kalan limit bilgisi payload'a eklenir.
\end{enumerate}

\subsection{Ask ve Blend Icin Farkli Maliyet Yaklasimi}
\begin{itemize}[leftmargin=*]
  \item \texttt{/api/ask}: \texttt{request\_weight=1}
  \item \texttt{/api/blend}: model sayisi kadar daha maliyetli kabul edilir
  (token tahmini model sayisi ile carpilir, agirlik en az 2 olacak sekilde artar)
\end{itemize}

Bu sayede blend gibi coklu-model akislarinda kota tuketimi gercege daha yakin
yansitilir.

\subsection{API Sozlesmesi (Frontend Için Anlasilir Ozet)}
\begin{itemize}[leftmargin=*]
  \item \texttt{GET /api/billing/usage}
  \begin{itemize}[leftmargin=*]
    \item plan, limitler, kullanilan ve kalan degerleri doner.
  \end{itemize}
  \item \texttt{POST /api/billing/plan}
  \begin{itemize}[leftmargin=*]
    \item \texttt{plan=free|premium} alir, plani gunceller, yeni usage ozetini doner.
  \end{itemize}
  \item Kota asiminda (ask/blend):
  \begin{itemize}[leftmargin=*]
    \item HTTP \texttt{429}
    \item \texttt{reason}: \texttt{daily\_requests\_exceeded} veya \texttt{monthly\_tokens\_exceeded}
    \item \texttt{upgrade\_required}: free planda \texttt{true}
  \end{itemize}
\end{itemize}

\subsection{Kullanici Deneyimi Tarafi}
\begin{itemize}[leftmargin=*]
  \item Tool drawer'da plan karti ve Free/Premium gecis butonlari bulunur.
  \item Kullanici plani degistirdiginde UI anlik olarak usage verisini yeniler.
  \item Limit asiminda kullaniciya teknik hata yerine yonlendirici mesaj verilir:
  "Daha fazlasi icin Premium'a gec" akisi.
\end{itemize}

\subsection{Kisa Senaryo (Once / Simdi)}
\begin{itemize}[leftmargin=*]
  \item \textbf{Once:} Her kullanici ayni akisla devam ediyor,
  yogun kullanimda maliyet kontrolu urun tarafinda net gorunmuyordu.
  \item \textbf{Simdi:} Free kullanici limitte 429 + yonlendirme mesaji aliyor,
  Premium kullanici daha yuksek limitlerle kesintisiz devam edebiliyor.
\end{itemize}

\subsection{Backend Kota Katmani}
\begin{itemize}[leftmargin=*]
  \item \texttt{server/app.py:1741} \texttt{consume\_plan\_quota(user, estimated\_tokens, request\_weight)}
  \item \texttt{server/app.py:1797} \texttt{GET /api/billing/usage}
  \item \texttt{server/app.py:1814} \texttt{POST /api/billing/plan}
  \item \texttt{server/app.py:3008} \texttt{/api/ask} kota kontrolu
  \item \texttt{server/app.py:3396} \texttt{/api/blend} kota kontrolu
\end{itemize}

\subsection{Frontend Plan Arayuzu}
\begin{itemize}[leftmargin=*]
  \item \texttt{client/src/App.jsx:216} \texttt{fetchUsageLimits}
  \item \texttt{client/src/App.jsx:231} \texttt{switchSubscriptionPlan}
  \item \texttt{client/src/App.jsx:2630-2636} Free/Premium gecis butonlari
\end{itemize}

\subsection{Davranis}
\begin{enumerate}[leftmargin=*]
  \item Kullanici her istek oncesi mevcut kota durumuna gore degerlendirilir.
  \item Limit asiminda istek reddedilir ve premiuma gecis mesaji donulur.
  \item Kullanici arac panelinden plan ve kullanim bilgisini gorebilir.
\end{enumerate}

\section{Kod Seviyesi Izlenebilirlik Matrisi}
\begin{longtable}{@{} P{3.5cm} P{4.9cm} P{6.2cm} @{} }
\toprule
Ozellik & Kod Noktasi & Sorumluluk \\
\midrule
Baglam optimizasyonu & \texttt{ChatInterface.jsx:194} & Prompt sikistirma ve rapor uretimi \\
Baglam optimizasyonu & \texttt{ChatInterface.jsx:1527} & Auto mode UI kontrolu \\
Patch risk skoru & \texttt{ChatInterface.jsx:241} & Risk analizi ve skor hesaplama \\
Patch guvenlik modal & \texttt{ChatInterface.jsx:2015} & Safe/Review/Block karar arayuzu \\
Canli review verisi & \texttt{models.py:72,85} & Review ve comment tablolari \\
Canli review API & \texttt{app.py:6195,6221,6256} & Getir/yorumla/durum degistir endpointleri \\
Plan kota kontrolu & \texttt{app.py:1741,3008,3396} & Ask/blend oncesi kota enforcement \\
Plan API + UI & \texttt{app.py:1797,1814} ve \texttt{App.jsx:216,231} & Kullanimi cekme ve plan degistirme \\
\bottomrule
\end{longtable}

\section{Test ve Dogrulama Sonuclari}
\subsection{Backend Smoke Test (Gercek API Cagrisi)}
Asagidaki kontroller, calisan backend uzerinden HTTP cagrilari ile dogrulanmistir.

\begin{longtable}{@{} P{4.4cm} P{3.0cm} P{6.8cm} @{} }
\toprule
Test Adimi & Sonuc & Kanit / Beklenen Davranis \\
\midrule
Kullanici kaydi ve JWT alma & PASS & \texttt{register\_ok=true} \\
Plan kullanim bilgisi & PASS & Baslangic plani \texttt{free} \\
Plan degistirme (free->premium->free) & PASS & \texttt{after\_premium=premium}, \texttt{after\_free=free} \\
Kota asimi (buyuk prompt) & PASS & HTTP \texttt{429}; \texttt{reason=monthly\_tokens\_exceeded}; \texttt{upgrade\_required=true} \\
Collaboration share token uretimi & PASS & \texttt{share\_token\_exists=true} \\
Review baslangic durumu & PASS & \texttt{review\_initial=open} \\
Review yorum ekleme & PASS & \texttt{comment\_author=smoke-guest} \\
Review durum guncelleme & PASS & \texttt{status\_after\_update=approved} \\
Review final kontrol & PASS & \texttt{review\_final=approved}, \texttt{comment\_count=1} \\
\bottomrule
\end{longtable}

\subsection{Frontend Manuel Test Plani (Kullanici Tarafi)}
Asagidaki adimlar UI davranisini birebir dogrulamak icin hazirlanmistir.

\textbf{Test-1: Baglam Optimizasyon Asistani}
\begin{enumerate}[leftmargin=*]
  \item Chat ekraninda \texttt{Auto Optimize} secenegini \texttt{ON} yapin.
  \item Tekrarli satirlar iceren uzun bir prompt yazin.
  \item Beklenen: optimize onerisi karti acilir, token azalimi raporlanir.
  \item Beklenen: \texttt{Apply Optimized Text}, \texttt{Send Optimized}, \texttt{Send Original} butonlari gorunur.
\end{enumerate}

\textbf{Test-2: Patch Guvenlik Seviyesi}
\begin{enumerate}[leftmargin=*]
  \item Kod degisikligi onerisi uretecek bir patch akisi tetikleyin.
  \item Beklenen: \texttt{Patch Safety Check} modal acilir.
  \item Beklenen: risk gerekceleri listelenir ve sonuca gore safe/review/block davranisi calisir.
\end{enumerate}

\textbf{Test-3: Canli Isbirligi Review Dongusu}
\begin{enumerate}[leftmargin=*]
  \item Bir konusma icin paylasim linki uretin.
  \item Paylasim ekranindan yorum ekleyin.
  \item Durumu sirayla \texttt{open -> revision\_requested -> approved} olarak degistirin.
  \item Beklenen: durum ve yorum listesi panelde guncel gorunur.
\end{enumerate}

\textbf{Test-4: Free/Premium Plan ve Kota}
\begin{enumerate}[leftmargin=*]
  \item Tool drawer uzerinden plani \texttt{free} ve \texttt{premium} arasinda degistirin.
  \item Beklenen: kullanim kartinda plan bilgisi anlik guncellenir.
  \item Free planda yuksek tokenli istek deneyin.
  \item Beklenen: limit asiminda premiuma yonlendiren hata mesaji gorunur.
\end{enumerate}

\section{MVP Sinirlari ve Sonraki Adimlar}
\begin{itemize}[leftmargin=*]
  \item Plan gecisi su an uygulama ici MVP toggledir; odeme altyapisi baglantisi sonraki fazdadir.
  \item Patch risk skoru heuristic tabanlidir; ileride AST ve test graph ile guclendirilebilir.
  \item Isbirligi paneli polling ile yenilenir; websocket/SSE ile daha dusuk gecikmeli hale getirilebilir.
\end{itemize}

\section{Sonuc}
8. hafta kapsami, 6. ve 7. haftada acilan alanlari eyleme donusturen bir ust katman olarak uygulanmistir:

\begin{itemize}[leftmargin=*]
  \item Context gorunurlugunden optimize edilebilir prompt akisina,
  \item Patch uygulamadan once guvenlik karar mekanizmasina,
  \item Tekil paylasimdan takim ici review dongusune,
  \item Sinirsiz kullanimdan plan/kota kontrollu urun modeline gecis saglanmistir.
\end{itemize}

\end{document}
